\documentclass[11pt,a4paper]{article}
\usepackage[utf8]{inputenc}
\usepackage[T1]{fontenc}
\usepackage[sc]{mathpazo}
\usepackage[scaled=0.92]{helvet}
\usepackage{microtype}
\usepackage[a4paper,top=2.6cm,bottom=2.4cm,left=2.4cm,right=2.4cm,headheight=14pt]{geometry}
\usepackage[table]{xcolor}
\usepackage{booktabs,array,enumitem,multirow}
\usepackage{tikz}
\usetikzlibrary{calc,positioning}
\usepackage{pgfplots}
\pgfplotsset{compat=1.16}
\usepackage[most]{tcolorbox}
\usepackage{titlesec}
\usepackage{needspace,etoolbox}
\usepackage{fancyhdr}
\usepackage{caption}
\usepackage{url}
\def\UrlBreaks{\do\/\do\-\do\_\do\.\do\=\do\?\do\&\do\a\do\b\do\c\do\d\do\e\do\f\do\g\do\h\do\i\do\j\do\k\do\l\do\m\do\n\do\o\do\p\do\q\do\r\do\s\do\t\do\u\do\v\do\w\do\x\do\y\do\z}
\Urlmuskip=0mu plus 1mu
\urlstyle{same}
\usepackage[hidelinks]{hyperref}

% ---- Palet (selaras dengan halaman dokumen pelengkap) ----
\definecolor{teal}{HTML}{0A6C6E}
\definecolor{tealdark}{HTML}{07484A}
\definecolor{ink}{HTML}{13262A}
\definecolor{muted}{HTML}{51666A}
\definecolor{rule}{HTML}{D3DEDC}
\definecolor{pale}{HTML}{EEF5F4}
\definecolor{amber}{HTML}{B7791F}
\definecolor{amberpale}{HTML}{FBF4E4}
\definecolor{well}{HTML}{CFE3E0}

\renewcommand{\tablename}{Tabel}
\renewcommand{\figurename}{Gambar}
\renewcommand{\refname}{Daftar Pustaka}
\color{ink}
\setlength{\parindent}{0pt}
\setlength{\parskip}{6.5pt plus 1pt}
\linespread{1.06}
\tolerance=1200
\emergencystretch=3em
\hyphenpenalty=1500
\setlist{itemsep=2.5pt,topsep=3pt,parsep=0pt,leftmargin=1.4em}
\setlist[itemize]{label={\color{teal}\textbullet}}
\renewcommand{\arraystretch}{1.22}
\setlength{\tabcolsep}{5pt}
\captionsetup{font={small,sf,color=muted},labelfont={bf,color=teal},labelsep=period,justification=raggedright,singlelinecheck=false,skip=5pt}

% ---- Judul bagian ----
\titleformat{\section}{\sffamily\bfseries\Large\color{ink}}{\color{teal}\thesection}{0.7em}{}[{\color{rule}\titlerule[0.6pt]}]
\titlespacing*{\section}{0pt}{18pt}{9pt}
\titleformat{\subsection}{\sffamily\bfseries\normalsize\color{tealdark}}{\color{teal}\thesubsection}{0.6em}{}
\titlespacing*{\subsection}{0pt}{11pt}{3pt}

% ---- Kepala dan kaki halaman ----
\pagestyle{fancy}
\fancyhf{}
\renewcommand{\headrulewidth}{0pt}
\fancyhead[L]{\sffamily\scriptsize\color{muted}\textsc{Alergen}\ \raisebox{0.15ex}{\tiny$\bullet$}\ Certificate of Analysis}
\fancyhead[R]{\sffamily\scriptsize\color{muted}Kit Diagnostik Alergi Berbasis Patch}
\fancyfoot[L]{\sffamily\scriptsize\color{muted}Universitas Airlangga \ $\cdot$ \ IDS000011522}
\fancyfoot[R]{\sffamily\small\bfseries\color{teal}\thepage}
\renewcommand{\headrule}{\color{rule}\hrule width\headwidth height 0.5pt}
\renewcommand{\headrulewidth}{0.5pt}
\fancypagestyle{plain}{\fancyhf{}\renewcommand{\headrulewidth}{0pt}\fancyfoot[R]{\sffamily\small\bfseries\color{teal}\thepage}}

% ---- Tabel ----
\newcommand{\tH}[1]{\textcolor{white}{\sffamily\bfseries\small #1}}
\newcommand{\tabstart}{\small\rowcolors{2}{white}{pale}\arrayrulecolor{rule}}
\newcommand{\hdrrow}{\rowcolor{tealdark}}
\newcolumntype{L}[1]{>{\raggedright\arraybackslash}p{#1}}
\newcolumntype{R}[1]{>{\raggedleft\arraybackslash}p{#1}}
\setlength{\aboverulesep}{0pt}\setlength{\belowrulesep}{0pt}
\renewcommand{\toprule}{\arrayrulecolor{tealdark}\specialrule{0.9pt}{0pt}{0pt}\arrayrulecolor{rule}}
\renewcommand{\bottomrule}{\arrayrulecolor{tealdark}\specialrule{0.9pt}{0pt}{0pt}\arrayrulecolor{rule}}

% ---- Kotak catatan dan kartu angka ----
\newtcolorbox{catatan}[1][Catatan]{enhanced,breakable,colback=amberpale,colframe=amberpale,boxrule=0pt,
  borderline west={2.2pt}{0pt}{amber},arc=1.5pt,left=10pt,right=8pt,top=6pt,bottom=6pt,
  before skip=8pt,after skip=10pt,
  title={\sffamily\bfseries\small\color{amber}\MakeUppercase{#1}},coltitle=amber,
  attach boxed title to top left={yshift=-2mm,xshift=0pt},boxed title style={empty,left=0pt,top=0pt,bottom=0pt},
  fonttitle=\sffamily,coltext=ink,before upper={\vspace{2pt}}}
\newtcolorbox{kartu}[2]{enhanced,colback=pale,colframe=pale,boxrule=0pt,arc=3pt,
  left=8pt,right=6pt,top=6pt,bottom=6pt,halign=left,
  borderline west={2.5pt}{0pt}{teal},
  before upper={{\sffamily\bfseries\LARGE\color{teal}#1}\par\vspace{1pt}},
  fontupper=\sffamily\scriptsize\color{muted},height=1.95cm,valign=top}
\newcommand{\kartuRow}[4]{%
 \begin{tcbraster}[raster columns=#1,raster equal height,raster column skip=7pt,raster row skip=7pt]#2#3#4\end{tcbraster}\vspace{9pt}}

% ---- Sampul ----
\newcommand{\patchmotif}[2]{% #1 skala, #2 warna garis
 \begin{scope}[scale=#1]
 \draw[#2,line width=1.2pt,rounded corners=9pt] (0,0) rectangle (15,6);
 \foreach \r in {0,1}{\foreach \c in {0,...,4}{
   \draw[#2,line width=1pt] ({1.9+\c*2.7},{4.1-\r*2.2}) circle (0.9);}}
 \end{scope}}

\newcommand{\makecover}[5]{% kode, judul, subjudul, tagline, ringkasan satu kalimat
\begin{titlepage}
\thispagestyle{empty}
\begin{tikzpicture}[remember picture,overlay]
 \fill[tealdark] (current page.north west) rectangle ($(current page.north east)+(0,-17.2cm)$);
 \fill[teal!85!tealdark] ($(current page.north east)+(0,-17.2cm)$) -- ($(current page.north east)+(-10cm,-17.2cm)$) -- ($(current page.north east)+(0,-10.8cm)$) -- cycle;
 \begin{scope}[shift={($(current page.north east)+(-9.9cm,-6.4cm)$)}]
   \patchmotif{0.5}{white!35}
 \end{scope}
 \node[anchor=north west,text=white!70,font=\sffamily\small\bfseries,align=left] at ($(current page.north west)+(2.4cm,-2.7cm)$) {\textsc{Alergen \ $\cdot$ \ Universitas Airlangga}};
 \node[anchor=north west,text=white,font=\sffamily\bfseries\fontsize{44}{44}\selectfont] at ($(current page.north west)+(2.4cm,-7.4cm)$) {#1};
 \node[anchor=north west,text=white,font=\sffamily\bfseries\fontsize{26}{31}\selectfont,text width=14.2cm,align=left] at ($(current page.north west)+(2.4cm,-10.2cm)$) {#2};
 \node[anchor=north west,text=white!85,font=\sffamily\fontsize{13}{17}\selectfont,text width=14cm,align=left] at ($(current page.north west)+(2.4cm,-14.1cm)$) {\hyphenpenalty=10000\exhyphenpenalty=10000 #3};
 \node[anchor=north west,text=ink,font=\fontsize{12}{17}\selectfont,text width=15.2cm,align=left] at ($(current page.north west)+(2.4cm,-18.8cm)$) {#4};
 \draw[teal,line width=2pt] ($(current page.north west)+(2.4cm,-21.9cm)$) -- ++(2.2cm,0);
 \node[anchor=north west,text=muted,font=\sffamily\small,text width=15cm,align=left] at ($(current page.north west)+(2.4cm,-22.4cm)$) {%
  \textbf{\color{ink}Kit Diagnostik Alergi Berbasis Patch (Alergen)}\\[3pt]
  Universitas Airlangga, Surabaya\\
  Paten Sederhana IDS000011522\\
  Oktober 2026};
 \node[anchor=south east,text=muted,font=\sffamily\scriptsize,align=right] at ($(current page.south east)+(-2.4cm,1.6cm)$) {#5};
\end{tikzpicture}
\end{titlepage}
\newpage}

\definecolor{sr1}{HTML}{8AA6AB}
\definecolor{sr2}{HTML}{3FA7A3}
\definecolor{sr3}{HTML}{0A6C6E}
\definecolor{sr4}{HTML}{B7791F}
\pgfplotsset{natbase/.style={width=0.95\linewidth,height=6.2cm,
  every axis label/.append style={font=\sffamily\scriptsize,color=muted},
  tick label style={font=\sffamily\scriptsize,color=muted},
  legend style={font=\sffamily\scriptsize,draw=rule,fill=white,at={(0.5,-0.2)},anchor=north},
  grid=major,grid style={color=rule,line width=0.4pt},
  axis line style={color=rule},tick style={color=rule},
  scaled ticks=false,/pgf/number format/.cd,1000 sep={.},fixed,precision=0,/tikz/.cd},
  natplot/.style={natbase,xtick={1,2,3,4,5},xticklabels={2028,2029,2030,2031,2032}}}
\hypersetup{pdftitle={Certificate of Analysis (CoA)},pdfauthor={Tim Alergen, Universitas Airlangga}}

\pretocmd{\section}{\Needspace{7\baselineskip}}{}{}
\pretocmd{\subsection}{\Needspace{4\baselineskip}}{}{}

\usepackage{float}
\fancyfoot[C]{\sffamily\tiny\color{muted}Dokumen contoh, nilai bukan hasil pengujian nyata}

\newcommand{\isi}[1][3.4cm]{\underline{\hspace{#1}}}
\newcommand{\coaHeader}[2]{% judul, subjudul
 \begin{tcolorbox}[enhanced,colback=tealdark,colframe=tealdark,arc=3pt,boxrule=0pt,left=12pt,right=12pt,top=9pt,bottom=9pt]
 {\sffamily\scriptsize\color{white!70}\textsc{Certificate of Analysis \ $\cdot$ \ Sertifikat Analisis}}\\[3pt]
 {\sffamily\bfseries\Large\color{white}#1}\\[2pt]
 {\sffamily\small\color{white!85}#2}
 \end{tcolorbox}}

\begin{document}

\makecover{CoA}{Certificate of Analysis (CoA)}{Sepuluh ekstrak alergen sebagai bahan baku dan patch Alergen sebagai produk jadi}{Dokumen ini memuat sertifikat analisis lengkap untuk sepuluh ekstrak alergen dan satu batch patch Alergen, beserta ringkasan batch, metode uji, dan kriteria penerimaan.}{Sepuluh CoA bahan baku, satu CoA produk jadi, dan lampiran hasil per sumur}

\section*{Ringkasan}
\addcontentsline{toc}{section}{Ringkasan}

\kartuRow{4}
{\begin{kartu}{10}{}ekstrak alergen bahan baku, masing-masing dengan satu CoA\end{kartu}}
{\begin{kartu}{1}{}batch patch Alergen sebagai produk jadi, 1.000 patch\end{kartu}}
{\begin{kartu}{30 $\mu$g}{}target protein per alergen per sumur, 10 sumur per patch\end{kartu}
\begin{kartu}{Lulus}{}seluruh 11 batch memenuhi spesifikasi\end{kartu}}

\emph{Certificate of Analysis} (CoA) adalah dokumen hasil pengujian satu batch tertentu yang ditandatangani bagian mutu. Dokumen ini berisi satu CoA untuk tiap ekstrak alergen yang menjadi bahan baku patch, satu CoA untuk produk jadi, dan lampiran hasil per sumur. Setiap CoA memuat nomor CoA dan nomor batch, tanggal produksi, uji, dan rilis, kedaluwarsa, kondisi penyimpanan, parameter dengan spesifikasi, hasil, metode, status lulus, kesimpulan, dan tanda tangan.



\clearpage
\section{Ringkasan batch}

Tabel~\ref{tab:ringkasan} mencantumkan batch yang dicakup, nomor CoA, kadar penanda alergen utama, aktivitas alergenik relatif terhadap ekstrak acuan internal (\emph{in-house reference preparation}, IHRP), dan tanggal rilis.

\begin{table}[H]
\centering
\caption{Ringkasan batch bahan baku dan produk jadi.}
\label{tab:ringkasan}
{\footnotesize\tabstart\setlength{\tabcolsep}{3.5pt}
\begin{tabular}{@{}rL{2.3cm}L{2.6cm}L{3.0cm}L{2.8cm}L{1.1cm}L{1.9cm}@{}}
\toprule
\hdrrow\tH{No} & \tH{Sediaan} & \tH{No. batch} & \tH{Penanda utama} & \tH{Aktivitas} & \tH{Status} & \tH{Rilis} \\
\midrule
1 & Tungau debu rumah & ALG-HDM-2608-01 & Der p 1 18,7 ng/mL & 97,4\% IHRP & Lulus & 14 September 2026 \\
2 & Udang & ALG-UDG-2608-01 & Tropomiosin 276,0 ng/mL & 102,8\% IHRP & Lulus & 15 September 2026 \\
3 & Susu sapi & ALG-SSP-2608-01 & Bos d 5 187,4 ng/mL & 99,1\% IHRP & Lulus & 16 September 2026 \\
4 & Daging sapi & ALG-DSP-2608-01 & Bos d 6 (albumin serum sapi) 342,0 ng/mL & 95,6\% IHRP & Lulus & 17 September 2026 \\
5 & Daging ayam & ALG-AYM-2608-01 & Gal d 5 (albumin serum ayam) 215,0 ng/mL & 101,3\% IHRP & Lulus & 18 September 2026 \\
6 & Putih telur & ALG-PTL-2608-01 & Ovalbumin (Gal d 2) 71,6 ng/mL & 98,2\% IHRP & Lulus & 19 September 2026 \\
7 & Gandum & ALG-GDM-2608-01 & Gliadin (Tri a 19) 168,0 ng/mL & 93,9\% IHRP & Lulus & 20 September 2026 \\
8 & Kacang tanah & ALG-KCT-2608-01 & Ara h 2 436,0 ng/mL & 104,5\% IHRP & Lulus & 21 September 2026 \\
9 & Kelapa & ALG-KLP-2608-01 & Coc n 1 126,0 ng/mL & 96,7\% IHRP & Lulus & 22 September 2026 \\
10 & Polen rumput bermuda & ALG-PLN-2608-01 & Cyn d 1 812,0 ng/mL & 100,9\% IHRP & Lulus & 23 September 2026 \\
11 & Patch Alergen & ALG-P-2610-01 & 10 penanda, Tabel~\ref{tab:sumur} & 94,6 sampai 103,7\% IHRP & Lulus & 22 Oktober 2026 \\
\bottomrule
\end{tabular}}
\end{table}

\clearpage
\section{Kriteria penerimaan dan struktur CoA}

Kriteria penerimaan tiap parameter ditetapkan sebelum batch diuji dan dicantumkan pada kolom spesifikasi. Struktur parameter mengikuti pedoman EMA untuk produk alergen \cite{ema}: identitas, kemurnian, aktivitas alergenik total, dan kadar alergen utama pada bahan aktif; kontrol mutu mikrobiologi pada produk. Tabel~\ref{tab:kriteria} merangkum kriteria yang berlaku seragam.

\begin{table}[H]
\centering
\caption{Kriteria penerimaan yang berlaku untuk seluruh ekstrak.}
\label{tab:kriteria}
{\footnotesize\tabstart\setlength{\tabcolsep}{4pt}
\begin{tabular}{@{}L{4.2cm}L{5.6cm}L{4.8cm}@{}}
\toprule
\hdrrow\tH{Parameter} & \tH{Kriteria} & \tH{Dasar} \\
\midrule
Volume per vial & 5,0 mL & Desain kemasan \\
Protein total & 0,5 sampai 2,5 mg/mL & Rentang kerja Bradford \\
Aktivitas alergenik & 80 sampai 125\% IHRP & Konsistensi antarbatch \cite{ema} \\
TAMC dan TYMC & $\leq 10^2$ dan $\leq 10^1$ CFU & Ph. Eur. 5.1.4, sediaan kulit \cite{ph514} \\
\emph{S. aureus}, \emph{P. aeruginosa} & Tidak ada dalam 1 g atau 1 mL & Ph. Eur. 5.1.4 \cite{ph514} \\
Endotoksin & $\leq 20$ EU/mL & Batas internal \\
Protein per sumur (patch) & 30 $\mu$g, 80 sampai 120\% & Desain produk \\
Keseragaman antarpatch & RSD $\leq 10\%$ & Batas internal \\
\bottomrule
\end{tabular}}
\end{table}

\begin{catatan}[Penanda alergen]
Penanda utama tiap ekstrak dipakai untuk identitas dan kadar: Der p 1 (tungau debu rumah), tropomiosin (udang), Bos d 5 (susu sapi), albumin serum sapi Bos d 6 (daging sapi), albumin serum ayam Gal d 5 (daging ayam), ovalbumin Gal d 2 (putih telur), gliadin Tri a 19 (gandum), Ara h 2 (kacang tanah), Coc n 1 (kelapa), dan Cyn d 1 (polen rumput bermuda).
\end{catatan}

\section{CoA bahan baku (ekstrak alergen)}

Sepuluh sertifikat berikut masing-masing berlaku untuk satu batch ekstrak, dengan volume 5 mL per vial dan pelarut dapar fosfat salin (PBS) pH 7,4. Ekstrak dibuat dari bahan sumber yang dikurangi lemaknya dengan aseton, diekstraksi dalam PBS, disaring, dan disimpan pada 2 sampai 8 $^\circ$C.

\clearpage
\subsection{Tungau debu rumah (\emph{House Dust Mite Extract})}

\coaHeader{House Dust Mite Extract}{Ekstrak alergen bahan baku patch Alergen \ $\cdot$ \ Tungau debu rumah}

\begin{tabular}{@{}l@{\hspace{8pt}}L{4.3cm}l@{\hspace{8pt}}L{4.3cm}@{}}
{\sffamily\small\color{muted}No. CoA} & {\sffamily\small\bfseries CoA/RM/HDM/01/IX/2026} & {\sffamily\small\color{muted}No. batch} & {\sffamily\small\bfseries ALG-HDM-2608-01} \\[2pt]
{\sffamily\small\color{muted}Tanggal produksi} & {\sffamily\small\bfseries 3 Agustus 2026} & {\sffamily\small\color{muted}Tanggal uji} & {\sffamily\small\bfseries 1 September 2026} \\[2pt]
{\sffamily\small\color{muted}Tanggal rilis} & {\sffamily\small\bfseries 14 September 2026} & {\sffamily\small\color{muted}Kedaluwarsa} & {\sffamily\small\bfseries Agustus 2027} \\[2pt]
{\sffamily\small\color{muted}Ukuran batch} & {\sffamily\small\bfseries 100 vial, 5 mL (500 mL)} & {\sffamily\small\color{muted}Penyimpanan} & {\sffamily\small\bfseries 2 sampai 8 $^\circ$C, terlindung cahaya} \\
\end{tabular}

\vspace{3pt}{\sffamily\small\color{muted}Sumber bahan: }{\small \emph{Dermatophagoides pteronyssinus}, tungau utuh dari biakan pada tepung jagung, panen minggu ke-3. Lot sumber HDM-S-01.}\par

\vspace{4pt}
{\footnotesize\tabstart\setlength{\tabcolsep}{3pt}\renewcommand{\arraystretch}{1.08}
\begin{tabular}{@{}rL{3.3cm}L{3.6cm}L{2.9cm}L{2.6cm}L{0.9cm}@{}}
\toprule

\hdrrow\tH{No} & \tH{Parameter} & \tH{Spesifikasi} & \tH{Hasil} & \tH{Metode} & \tH{Status} \\
\midrule

\rowcolor{rule!55}\multicolumn{6}{@{}l}{\sffamily\scriptsize\bfseries\color{tealdark}\MakeUppercase{A. Sifat fisik dan identitas}} \\
1 & Bentuk dan warna & Jernih, tidak berwarna & Jernih, tidak berwarna & Visual & Lulus \\
2 & Volume per vial & 5,0 mL & 5,0 mL & Volumetrik & Lulus \\
3 & Kelarutan & Larut dalam air & Larut sempurna & Visual & Lulus \\
4 & Keasaman (pH) & 7,4 sampai 7,8 & 7,6 & Potensiometri & Lulus \\
\rowcolor{rule!55}\multicolumn{6}{@{}l}{\sffamily\scriptsize\bfseries\color{tealdark}\MakeUppercase{B. Karakterisasi kimia dan biologi}} \\
5 & Nitrogen total & 0,08 sampai 0,11\% & 0,094\% & Kjeldahl & Lulus \\
6 & Protein total & 0,5 sampai 2,5 mg/mL & 1,12 mg/mL & Bradford \cite{ema} & Lulus \\
7 & Profil protein & Pita utama 25 kDa & Sesuai, 25 kDa & SDS-PAGE 12\% & Lulus \\
8 & Kadar Der p 1 & 11,3 sampai 26,6 ng/mL & 18,7 ng/mL & ELISA Der p 1 \cite{pratama23} & Lulus \\
9 & Titer IgE dan IgG4 & $>20$ IU/mL & 34 IU/mL & ELISA serum pool & Lulus \\
10 & Aktivitas alergenik & 80 sampai 125\% IHRP & 97,4\% IHRP & IgE-inhibisi \cite{ema} & Lulus \\
\rowcolor{rule!55}\multicolumn{6}{@{}l}{\sffamily\scriptsize\bfseries\color{tealdark}\MakeUppercase{C. Mutu mikrobiologi dan endotoksin}} \\
11 & Angka lempeng aerob (TAMC) & $\leq 10^2$ CFU/mL & $<10$ CFU/mL & Ph. Eur. 2.6.12 & Lulus \\
12 & Kapang dan khamir (TYMC) & $\leq 10^1$ CFU/mL & $<10$ CFU/mL & Ph. Eur. 2.6.12 & Lulus \\
13 & \emph{Staphylococcus aureus} & Tidak ada per 1 mL & Tidak ada & Ph. Eur. 2.6.13 & Lulus \\
14 & \emph{Pseudomonas aeruginosa} & Tidak ada per 1 mL & Tidak ada & Ph. Eur. 2.6.13 & Lulus \\
15 & Endotoksin & $\leq 20$ EU/mL & 2,1 EU/mL & Ph. Eur. 2.6.14 & Lulus \\

\bottomrule
\end{tabular}}

\vspace{6pt}

\textbf{Kesimpulan.} Batch ALG-HDM-2608-01 memenuhi seluruh spesifikasi dan dirilis sebagai bahan baku patch Alergen.

\vspace{4pt}
\begin{tabular}{@{}p{5.1cm}p{5.1cm}p{5.1cm}@{}}
\rule{0pt}{16pt}\isi[4.5cm] & \isi[4.5cm] & \isi[4.5cm]\\
{\sffamily\scriptsize\color{muted}Analis pengendalian mutu} & {\sffamily\scriptsize\color{muted}Kepala bagian mutu} & {\sffamily\scriptsize\color{muted}Penanggung jawab produksi}\\
{\sffamily\scriptsize\color{muted}Tanggal: 14 September 2026} & {\sffamily\scriptsize\color{muted}Tanggal: 14 September 2026} & {\sffamily\scriptsize\color{muted}Tanggal: 14 September 2026}
\end{tabular}


\clearpage
\subsection{Udang (\emph{Shrimp Extract})}

\coaHeader{Shrimp Extract}{Ekstrak alergen bahan baku patch Alergen \ $\cdot$ \ Udang}

\begin{tabular}{@{}l@{\hspace{8pt}}L{4.3cm}l@{\hspace{8pt}}L{4.3cm}@{}}
{\sffamily\small\color{muted}No. CoA} & {\sffamily\small\bfseries CoA/RM/UDG/02/IX/2026} & {\sffamily\small\color{muted}No. batch} & {\sffamily\small\bfseries ALG-UDG-2608-01} \\[2pt]
{\sffamily\small\color{muted}Tanggal produksi} & {\sffamily\small\bfseries 4 Agustus 2026} & {\sffamily\small\color{muted}Tanggal uji} & {\sffamily\small\bfseries 2 September 2026} \\[2pt]
{\sffamily\small\color{muted}Tanggal rilis} & {\sffamily\small\bfseries 15 September 2026} & {\sffamily\small\color{muted}Kedaluwarsa} & {\sffamily\small\bfseries Agustus 2027} \\[2pt]
{\sffamily\small\color{muted}Ukuran batch} & {\sffamily\small\bfseries 100 vial, 5 mL (500 mL)} & {\sffamily\small\color{muted}Penyimpanan} & {\sffamily\small\bfseries 2 sampai 8 $^\circ$C, terlindung cahaya} \\
\end{tabular}

\vspace{3pt}{\sffamily\small\color{muted}Sumber bahan: }{\small \emph{Penaeus vannamei}, daging udang umur 90 sampai 120 hari, sekitar 500 g. Lot sumber UDG-S-02.}\par

\vspace{4pt}
{\footnotesize\tabstart\setlength{\tabcolsep}{3pt}\renewcommand{\arraystretch}{1.08}
\begin{tabular}{@{}rL{3.3cm}L{3.6cm}L{2.9cm}L{2.6cm}L{0.9cm}@{}}
\toprule

\hdrrow\tH{No} & \tH{Parameter} & \tH{Spesifikasi} & \tH{Hasil} & \tH{Metode} & \tH{Status} \\
\midrule

\rowcolor{rule!55}\multicolumn{6}{@{}l}{\sffamily\scriptsize\bfseries\color{tealdark}\MakeUppercase{A. Sifat fisik dan identitas}} \\
1 & Bentuk dan warna & Jernih, tidak berwarna & Jernih, tidak berwarna & Visual & Lulus \\
2 & Volume per vial & 5,0 mL & 5,0 mL & Volumetrik & Lulus \\
3 & Kelarutan & Larut dalam air & Larut sempurna & Visual & Lulus \\
4 & Keasaman (pH) & 7,2 sampai 7,6 & 7,4 & Potensiometri & Lulus \\
\rowcolor{rule!55}\multicolumn{6}{@{}l}{\sffamily\scriptsize\bfseries\color{tealdark}\MakeUppercase{B. Karakterisasi kimia dan biologi}} \\
5 & Nitrogen total & 0,02 sampai 0,06\% & 0,041\% & Kjeldahl & Lulus \\
6 & Protein total & 0,5 sampai 2,5 mg/mL & 0,86 mg/mL & Bradford \cite{ema} & Lulus \\
7 & Profil protein & Pita utama 36 kDa & Sesuai, 36 kDa & SDS-PAGE 12\% & Lulus \\
8 & Kadar Tropomiosin & 100,0 sampai 500,0 ng/mL & 276,0 ng/mL & ELISA tropomiosin \cite{marhaeny24} & Lulus \\
9 & Aktivitas alergenik & 80 sampai 125\% IHRP & 102,8\% IHRP & IgE-inhibisi \cite{ema} & Lulus \\
\rowcolor{rule!55}\multicolumn{6}{@{}l}{\sffamily\scriptsize\bfseries\color{tealdark}\MakeUppercase{C. Mutu mikrobiologi dan endotoksin}} \\
10 & Angka lempeng aerob (TAMC) & $\leq 10^2$ CFU/mL & $<10$ CFU/mL & Ph. Eur. 2.6.12 & Lulus \\
11 & Kapang dan khamir (TYMC) & $\leq 10^1$ CFU/mL & $<10$ CFU/mL & Ph. Eur. 2.6.12 & Lulus \\
12 & \emph{Staphylococcus aureus} & Tidak ada per 1 mL & Tidak ada & Ph. Eur. 2.6.13 & Lulus \\
13 & \emph{Pseudomonas aeruginosa} & Tidak ada per 1 mL & Tidak ada & Ph. Eur. 2.6.13 & Lulus \\
14 & Endotoksin & $\leq 20$ EU/mL & 2,4 EU/mL & Ph. Eur. 2.6.14 & Lulus \\

\bottomrule
\end{tabular}}

\vspace{6pt}

\textbf{Kesimpulan.} Batch ALG-UDG-2608-01 memenuhi seluruh spesifikasi dan dirilis sebagai bahan baku patch Alergen.

\vspace{4pt}
\begin{tabular}{@{}p{5.1cm}p{5.1cm}p{5.1cm}@{}}
\rule{0pt}{16pt}\isi[4.5cm] & \isi[4.5cm] & \isi[4.5cm]\\
{\sffamily\scriptsize\color{muted}Analis pengendalian mutu} & {\sffamily\scriptsize\color{muted}Kepala bagian mutu} & {\sffamily\scriptsize\color{muted}Penanggung jawab produksi}\\
{\sffamily\scriptsize\color{muted}Tanggal: 15 September 2026} & {\sffamily\scriptsize\color{muted}Tanggal: 15 September 2026} & {\sffamily\scriptsize\color{muted}Tanggal: 15 September 2026}
\end{tabular}


\clearpage
\subsection{Susu sapi (\emph{Cow's Milk Extract})}

\coaHeader{Cow's Milk Extract}{Ekstrak alergen bahan baku patch Alergen \ $\cdot$ \ Susu sapi}

\begin{tabular}{@{}l@{\hspace{8pt}}L{4.3cm}l@{\hspace{8pt}}L{4.3cm}@{}}
{\sffamily\small\color{muted}No. CoA} & {\sffamily\small\bfseries CoA/RM/SSP/03/IX/2026} & {\sffamily\small\color{muted}No. batch} & {\sffamily\small\bfseries ALG-SSP-2608-01} \\[2pt]
{\sffamily\small\color{muted}Tanggal produksi} & {\sffamily\small\bfseries 5 Agustus 2026} & {\sffamily\small\color{muted}Tanggal uji} & {\sffamily\small\bfseries 3 September 2026} \\[2pt]
{\sffamily\small\color{muted}Tanggal rilis} & {\sffamily\small\bfseries 16 September 2026} & {\sffamily\small\color{muted}Kedaluwarsa} & {\sffamily\small\bfseries Agustus 2027} \\[2pt]
{\sffamily\small\color{muted}Ukuran batch} & {\sffamily\small\bfseries 100 vial, 5 mL (500 mL)} & {\sffamily\small\color{muted}Penyimpanan} & {\sffamily\small\bfseries 2 sampai 8 $^\circ$C, terlindung cahaya} \\
\end{tabular}

\vspace{3pt}{\sffamily\small\color{muted}Sumber bahan: }{\small Susu sapi segar, \emph{Bos taurus}. Lot sumber SSP-S-03.}\par

\vspace{4pt}
{\footnotesize\tabstart\setlength{\tabcolsep}{3pt}\renewcommand{\arraystretch}{1.08}
\begin{tabular}{@{}rL{3.3cm}L{3.6cm}L{2.9cm}L{2.6cm}L{0.9cm}@{}}
\toprule

\hdrrow\tH{No} & \tH{Parameter} & \tH{Spesifikasi} & \tH{Hasil} & \tH{Metode} & \tH{Status} \\
\midrule

\rowcolor{rule!55}\multicolumn{6}{@{}l}{\sffamily\scriptsize\bfseries\color{tealdark}\MakeUppercase{A. Sifat fisik dan identitas}} \\
1 & Bentuk dan warna & Jernih, tidak berwarna & Jernih, tidak berwarna & Visual & Lulus \\
2 & Volume per vial & 5,0 mL & 5,0 mL & Volumetrik & Lulus \\
3 & Kelarutan & Larut dalam air & Larut sempurna & Visual & Lulus \\
4 & Keasaman (pH) & 6,0 sampai 6,5 & 6,3 & Potensiometri & Lulus \\
\rowcolor{rule!55}\multicolumn{6}{@{}l}{\sffamily\scriptsize\bfseries\color{tealdark}\MakeUppercase{B. Karakterisasi kimia dan biologi}} \\
5 & Nitrogen total & 0,05 sampai 0,06\% & 0,055\% & Kjeldahl & Lulus \\
6 & Protein total & 0,5 sampai 2,5 mg/mL & 1,45 mg/mL & Bradford \cite{ema} & Lulus \\
7 & Profil protein & Pita utama 18 kDa & Sesuai, 18 kDa & SDS-PAGE 12\% & Lulus \\
8 & Kadar Bos d 5 & 42,6 sampai 331,8 ng/mL & 187,4 ng/mL & ELISA Bos d 5 (sandwich) & Lulus \\
9 & Aktivitas alergenik & 80 sampai 125\% IHRP & 99,1\% IHRP & IgE-inhibisi \cite{ema} & Lulus \\
\rowcolor{rule!55}\multicolumn{6}{@{}l}{\sffamily\scriptsize\bfseries\color{tealdark}\MakeUppercase{C. Mutu mikrobiologi dan endotoksin}} \\
10 & Angka lempeng aerob (TAMC) & $\leq 10^2$ CFU/mL & $<10$ CFU/mL & Ph. Eur. 2.6.12 & Lulus \\
11 & Kapang dan khamir (TYMC) & $\leq 10^1$ CFU/mL & $<10$ CFU/mL & Ph. Eur. 2.6.12 & Lulus \\
12 & \emph{Staphylococcus aureus} & Tidak ada per 1 mL & Tidak ada & Ph. Eur. 2.6.13 & Lulus \\
13 & \emph{Pseudomonas aeruginosa} & Tidak ada per 1 mL & Tidak ada & Ph. Eur. 2.6.13 & Lulus \\
14 & Endotoksin & $\leq 20$ EU/mL & 2,7 EU/mL & Ph. Eur. 2.6.14 & Lulus \\

\bottomrule
\end{tabular}}

\vspace{6pt}

\textbf{Kesimpulan.} Batch ALG-SSP-2608-01 memenuhi seluruh spesifikasi dan dirilis sebagai bahan baku patch Alergen.

\vspace{4pt}
\begin{tabular}{@{}p{5.1cm}p{5.1cm}p{5.1cm}@{}}
\rule{0pt}{16pt}\isi[4.5cm] & \isi[4.5cm] & \isi[4.5cm]\\
{\sffamily\scriptsize\color{muted}Analis pengendalian mutu} & {\sffamily\scriptsize\color{muted}Kepala bagian mutu} & {\sffamily\scriptsize\color{muted}Penanggung jawab produksi}\\
{\sffamily\scriptsize\color{muted}Tanggal: 16 September 2026} & {\sffamily\scriptsize\color{muted}Tanggal: 16 September 2026} & {\sffamily\scriptsize\color{muted}Tanggal: 16 September 2026}
\end{tabular}


\clearpage
\subsection{Daging sapi (\emph{Beef Extract})}

\coaHeader{Beef Extract}{Ekstrak alergen bahan baku patch Alergen \ $\cdot$ \ Daging sapi}

\begin{tabular}{@{}l@{\hspace{8pt}}L{4.3cm}l@{\hspace{8pt}}L{4.3cm}@{}}
{\sffamily\small\color{muted}No. CoA} & {\sffamily\small\bfseries CoA/RM/DSP/04/IX/2026} & {\sffamily\small\color{muted}No. batch} & {\sffamily\small\bfseries ALG-DSP-2608-01} \\[2pt]
{\sffamily\small\color{muted}Tanggal produksi} & {\sffamily\small\bfseries 6 Agustus 2026} & {\sffamily\small\color{muted}Tanggal uji} & {\sffamily\small\bfseries 4 September 2026} \\[2pt]
{\sffamily\small\color{muted}Tanggal rilis} & {\sffamily\small\bfseries 17 September 2026} & {\sffamily\small\color{muted}Kedaluwarsa} & {\sffamily\small\bfseries Agustus 2027} \\[2pt]
{\sffamily\small\color{muted}Ukuran batch} & {\sffamily\small\bfseries 100 vial, 5 mL (500 mL)} & {\sffamily\small\color{muted}Penyimpanan} & {\sffamily\small\bfseries 2 sampai 8 $^\circ$C, terlindung cahaya} \\
\end{tabular}

\vspace{3pt}{\sffamily\small\color{muted}Sumber bahan: }{\small Otot rangka sapi, \emph{Bos taurus}, bagian has dalam. Lot sumber DSP-S-04.}\par

\vspace{4pt}
{\footnotesize\tabstart\setlength{\tabcolsep}{3pt}\renewcommand{\arraystretch}{1.08}
\begin{tabular}{@{}rL{3.3cm}L{3.6cm}L{2.9cm}L{2.6cm}L{0.9cm}@{}}
\toprule

\hdrrow\tH{No} & \tH{Parameter} & \tH{Spesifikasi} & \tH{Hasil} & \tH{Metode} & \tH{Status} \\
\midrule

\rowcolor{rule!55}\multicolumn{6}{@{}l}{\sffamily\scriptsize\bfseries\color{tealdark}\MakeUppercase{A. Sifat fisik dan identitas}} \\
1 & Bentuk dan warna & Jernih, kemerahan muda & Jernih, kemerahan muda & Visual & Lulus \\
2 & Volume per vial & 5,0 mL & 5,0 mL & Volumetrik & Lulus \\
3 & Kelarutan & Larut dalam air & Larut sempurna & Visual & Lulus \\
4 & Keasaman (pH) & 7,0 sampai 7,6 & 7,3 & Potensiometri & Lulus \\
\rowcolor{rule!55}\multicolumn{6}{@{}l}{\sffamily\scriptsize\bfseries\color{tealdark}\MakeUppercase{B. Karakterisasi kimia dan biologi}} \\
5 & Nitrogen total & 0,03 sampai 0,07\% & 0,052\% & Kjeldahl & Lulus \\
6 & Protein total & 0,5 sampai 2,5 mg/mL & 1,02 mg/mL & Bradford \cite{ema} & Lulus \\
7 & Profil protein & Pita utama 67 kDa & Sesuai, 67 kDa & SDS-PAGE 12\% & Lulus \\
8 & Kadar Bos d 6 (albumin serum sapi) & 100,0 sampai 600,0 ng/mL & 342,0 ng/mL & ELISA Bos d 6 (sandwich) & Lulus \\
9 & Aktivitas alergenik & 80 sampai 125\% IHRP & 95,6\% IHRP & IgE-inhibisi \cite{ema} & Lulus \\
\rowcolor{rule!55}\multicolumn{6}{@{}l}{\sffamily\scriptsize\bfseries\color{tealdark}\MakeUppercase{C. Mutu mikrobiologi dan endotoksin}} \\
10 & Angka lempeng aerob (TAMC) & $\leq 10^2$ CFU/mL & $<10$ CFU/mL & Ph. Eur. 2.6.12 & Lulus \\
11 & Kapang dan khamir (TYMC) & $\leq 10^1$ CFU/mL & $<10$ CFU/mL & Ph. Eur. 2.6.12 & Lulus \\
12 & \emph{Staphylococcus aureus} & Tidak ada per 1 mL & Tidak ada & Ph. Eur. 2.6.13 & Lulus \\
13 & \emph{Pseudomonas aeruginosa} & Tidak ada per 1 mL & Tidak ada & Ph. Eur. 2.6.13 & Lulus \\
14 & Endotoksin & $\leq 20$ EU/mL & 3,0 EU/mL & Ph. Eur. 2.6.14 & Lulus \\

\bottomrule
\end{tabular}}

\vspace{6pt}

\textbf{Kesimpulan.} Batch ALG-DSP-2608-01 memenuhi seluruh spesifikasi dan dirilis sebagai bahan baku patch Alergen.

\vspace{4pt}
\begin{tabular}{@{}p{5.1cm}p{5.1cm}p{5.1cm}@{}}
\rule{0pt}{16pt}\isi[4.5cm] & \isi[4.5cm] & \isi[4.5cm]\\
{\sffamily\scriptsize\color{muted}Analis pengendalian mutu} & {\sffamily\scriptsize\color{muted}Kepala bagian mutu} & {\sffamily\scriptsize\color{muted}Penanggung jawab produksi}\\
{\sffamily\scriptsize\color{muted}Tanggal: 17 September 2026} & {\sffamily\scriptsize\color{muted}Tanggal: 17 September 2026} & {\sffamily\scriptsize\color{muted}Tanggal: 17 September 2026}
\end{tabular}


\clearpage
\subsection{Daging ayam (\emph{Chicken Meat Extract})}

\coaHeader{Chicken Meat Extract}{Ekstrak alergen bahan baku patch Alergen \ $\cdot$ \ Daging ayam}

\begin{tabular}{@{}l@{\hspace{8pt}}L{4.3cm}l@{\hspace{8pt}}L{4.3cm}@{}}
{\sffamily\small\color{muted}No. CoA} & {\sffamily\small\bfseries CoA/RM/AYM/05/IX/2026} & {\sffamily\small\color{muted}No. batch} & {\sffamily\small\bfseries ALG-AYM-2608-01} \\[2pt]
{\sffamily\small\color{muted}Tanggal produksi} & {\sffamily\small\bfseries 7 Agustus 2026} & {\sffamily\small\color{muted}Tanggal uji} & {\sffamily\small\bfseries 5 September 2026} \\[2pt]
{\sffamily\small\color{muted}Tanggal rilis} & {\sffamily\small\bfseries 18 September 2026} & {\sffamily\small\color{muted}Kedaluwarsa} & {\sffamily\small\bfseries Agustus 2027} \\[2pt]
{\sffamily\small\color{muted}Ukuran batch} & {\sffamily\small\bfseries 100 vial, 5 mL (500 mL)} & {\sffamily\small\color{muted}Penyimpanan} & {\sffamily\small\bfseries 2 sampai 8 $^\circ$C, terlindung cahaya} \\
\end{tabular}

\vspace{3pt}{\sffamily\small\color{muted}Sumber bahan: }{\small Otot dada ayam, \emph{Gallus gallus domesticus}. Lot sumber AYM-S-05.}\par

\vspace{4pt}
{\footnotesize\tabstart\setlength{\tabcolsep}{3pt}\renewcommand{\arraystretch}{1.08}
\begin{tabular}{@{}rL{3.3cm}L{3.6cm}L{2.9cm}L{2.6cm}L{0.9cm}@{}}
\toprule

\hdrrow\tH{No} & \tH{Parameter} & \tH{Spesifikasi} & \tH{Hasil} & \tH{Metode} & \tH{Status} \\
\midrule

\rowcolor{rule!55}\multicolumn{6}{@{}l}{\sffamily\scriptsize\bfseries\color{tealdark}\MakeUppercase{A. Sifat fisik dan identitas}} \\
1 & Bentuk dan warna & Jernih, tidak berwarna & Jernih, tidak berwarna & Visual & Lulus \\
2 & Volume per vial & 5,0 mL & 5,0 mL & Volumetrik & Lulus \\
3 & Kelarutan & Larut dalam air & Larut sempurna & Visual & Lulus \\
4 & Keasaman (pH) & 7,0 sampai 7,6 & 7,2 & Potensiometri & Lulus \\
\rowcolor{rule!55}\multicolumn{6}{@{}l}{\sffamily\scriptsize\bfseries\color{tealdark}\MakeUppercase{B. Karakterisasi kimia dan biologi}} \\
5 & Nitrogen total & 0,03 sampai 0,07\% & 0,048\% & Kjeldahl & Lulus \\
6 & Protein total & 0,5 sampai 2,5 mg/mL & 0,94 mg/mL & Bradford \cite{ema} & Lulus \\
7 & Profil protein & Pita utama 69 kDa & Sesuai, 69 kDa & SDS-PAGE 12\% & Lulus \\
8 & Kadar Gal d 5 (albumin serum ayam) & 80,0 sampai 500,0 ng/mL & 215,0 ng/mL & ELISA Gal d 5 (sandwich) & Lulus \\
9 & Aktivitas alergenik & 80 sampai 125\% IHRP & 101,3\% IHRP & IgE-inhibisi \cite{ema} & Lulus \\
\rowcolor{rule!55}\multicolumn{6}{@{}l}{\sffamily\scriptsize\bfseries\color{tealdark}\MakeUppercase{C. Mutu mikrobiologi dan endotoksin}} \\
10 & Angka lempeng aerob (TAMC) & $\leq 10^2$ CFU/mL & $<10$ CFU/mL & Ph. Eur. 2.6.12 & Lulus \\
11 & Kapang dan khamir (TYMC) & $\leq 10^1$ CFU/mL & $<10$ CFU/mL & Ph. Eur. 2.6.12 & Lulus \\
12 & \emph{Staphylococcus aureus} & Tidak ada per 1 mL & Tidak ada & Ph. Eur. 2.6.13 & Lulus \\
13 & \emph{Pseudomonas aeruginosa} & Tidak ada per 1 mL & Tidak ada & Ph. Eur. 2.6.13 & Lulus \\
14 & Endotoksin & $\leq 20$ EU/mL & 3,3 EU/mL & Ph. Eur. 2.6.14 & Lulus \\

\bottomrule
\end{tabular}}

\vspace{6pt}

\textbf{Kesimpulan.} Batch ALG-AYM-2608-01 memenuhi seluruh spesifikasi dan dirilis sebagai bahan baku patch Alergen.

\vspace{4pt}
\begin{tabular}{@{}p{5.1cm}p{5.1cm}p{5.1cm}@{}}
\rule{0pt}{16pt}\isi[4.5cm] & \isi[4.5cm] & \isi[4.5cm]\\
{\sffamily\scriptsize\color{muted}Analis pengendalian mutu} & {\sffamily\scriptsize\color{muted}Kepala bagian mutu} & {\sffamily\scriptsize\color{muted}Penanggung jawab produksi}\\
{\sffamily\scriptsize\color{muted}Tanggal: 18 September 2026} & {\sffamily\scriptsize\color{muted}Tanggal: 18 September 2026} & {\sffamily\scriptsize\color{muted}Tanggal: 18 September 2026}
\end{tabular}


\clearpage
\subsection{Putih telur (\emph{Egg White Extract})}

\coaHeader{Egg White Extract}{Ekstrak alergen bahan baku patch Alergen \ $\cdot$ \ Putih telur}

\begin{tabular}{@{}l@{\hspace{8pt}}L{4.3cm}l@{\hspace{8pt}}L{4.3cm}@{}}
{\sffamily\small\color{muted}No. CoA} & {\sffamily\small\bfseries CoA/RM/PTL/06/IX/2026} & {\sffamily\small\color{muted}No. batch} & {\sffamily\small\bfseries ALG-PTL-2608-01} \\[2pt]
{\sffamily\small\color{muted}Tanggal produksi} & {\sffamily\small\bfseries 8 Agustus 2026} & {\sffamily\small\color{muted}Tanggal uji} & {\sffamily\small\bfseries 6 September 2026} \\[2pt]
{\sffamily\small\color{muted}Tanggal rilis} & {\sffamily\small\bfseries 19 September 2026} & {\sffamily\small\color{muted}Kedaluwarsa} & {\sffamily\small\bfseries Agustus 2027} \\[2pt]
{\sffamily\small\color{muted}Ukuran batch} & {\sffamily\small\bfseries 100 vial, 5 mL (500 mL)} & {\sffamily\small\color{muted}Penyimpanan} & {\sffamily\small\bfseries 2 sampai 8 $^\circ$C, terlindung cahaya} \\
\end{tabular}

\vspace{3pt}{\sffamily\small\color{muted}Sumber bahan: }{\small Putih telur ayam, \emph{Gallus gallus domesticus}. Lot sumber PTL-S-06.}\par

\vspace{4pt}
{\footnotesize\tabstart\setlength{\tabcolsep}{3pt}\renewcommand{\arraystretch}{1.08}
\begin{tabular}{@{}rL{3.3cm}L{3.6cm}L{2.9cm}L{2.6cm}L{0.9cm}@{}}
\toprule

\hdrrow\tH{No} & \tH{Parameter} & \tH{Spesifikasi} & \tH{Hasil} & \tH{Metode} & \tH{Status} \\
\midrule

\rowcolor{rule!55}\multicolumn{6}{@{}l}{\sffamily\scriptsize\bfseries\color{tealdark}\MakeUppercase{A. Sifat fisik dan identitas}} \\
1 & Bentuk dan warna & Jernih, tidak berwarna & Jernih, tidak berwarna & Visual & Lulus \\
2 & Volume per vial & 5,0 mL & 5,0 mL & Volumetrik & Lulus \\
3 & Kelarutan & Larut dalam air & Larut sempurna & Visual & Lulus \\
4 & Keasaman (pH) & 7,8 sampai 8,0 & 7,9 & Potensiometri & Lulus \\
\rowcolor{rule!55}\multicolumn{6}{@{}l}{\sffamily\scriptsize\bfseries\color{tealdark}\MakeUppercase{B. Karakterisasi kimia dan biologi}} \\
5 & Nitrogen total & 0,02 sampai 0,05\% & 0,034\% & Kjeldahl & Lulus \\
6 & Protein total & 0,5 sampai 2,5 mg/mL & 0,78 mg/mL & Bradford \cite{ema} & Lulus \\
7 & Profil protein & Pita utama 45 kDa & Sesuai, 45 kDa & SDS-PAGE 12\% & Lulus \\
8 & Kadar Ovalbumin (Gal d 2) & 50,0 sampai 100,0 ng/mL & 71,6 ng/mL & ELISA ovalbumin (sandwich) & Lulus \\
9 & Aktivitas alergenik & 80 sampai 125\% IHRP & 98,2\% IHRP & IgE-inhibisi \cite{ema} & Lulus \\
\rowcolor{rule!55}\multicolumn{6}{@{}l}{\sffamily\scriptsize\bfseries\color{tealdark}\MakeUppercase{C. Mutu mikrobiologi dan endotoksin}} \\
10 & Angka lempeng aerob (TAMC) & $\leq 10^2$ CFU/mL & $<10$ CFU/mL & Ph. Eur. 2.6.12 & Lulus \\
11 & Kapang dan khamir (TYMC) & $\leq 10^1$ CFU/mL & $<10$ CFU/mL & Ph. Eur. 2.6.12 & Lulus \\
12 & \emph{Staphylococcus aureus} & Tidak ada per 1 mL & Tidak ada & Ph. Eur. 2.6.13 & Lulus \\
13 & \emph{Pseudomonas aeruginosa} & Tidak ada per 1 mL & Tidak ada & Ph. Eur. 2.6.13 & Lulus \\
14 & Endotoksin & $\leq 20$ EU/mL & 2,1 EU/mL & Ph. Eur. 2.6.14 & Lulus \\

\bottomrule
\end{tabular}}

\vspace{6pt}

\textbf{Kesimpulan.} Batch ALG-PTL-2608-01 memenuhi seluruh spesifikasi dan dirilis sebagai bahan baku patch Alergen.

\vspace{4pt}
\begin{tabular}{@{}p{5.1cm}p{5.1cm}p{5.1cm}@{}}
\rule{0pt}{16pt}\isi[4.5cm] & \isi[4.5cm] & \isi[4.5cm]\\
{\sffamily\scriptsize\color{muted}Analis pengendalian mutu} & {\sffamily\scriptsize\color{muted}Kepala bagian mutu} & {\sffamily\scriptsize\color{muted}Penanggung jawab produksi}\\
{\sffamily\scriptsize\color{muted}Tanggal: 19 September 2026} & {\sffamily\scriptsize\color{muted}Tanggal: 19 September 2026} & {\sffamily\scriptsize\color{muted}Tanggal: 19 September 2026}
\end{tabular}


\clearpage
\subsection{Gandum (\emph{Wheat Extract})}

\coaHeader{Wheat Extract}{Ekstrak alergen bahan baku patch Alergen \ $\cdot$ \ Gandum}

\begin{tabular}{@{}l@{\hspace{8pt}}L{4.3cm}l@{\hspace{8pt}}L{4.3cm}@{}}
{\sffamily\small\color{muted}No. CoA} & {\sffamily\small\bfseries CoA/RM/GDM/07/IX/2026} & {\sffamily\small\color{muted}No. batch} & {\sffamily\small\bfseries ALG-GDM-2608-01} \\[2pt]
{\sffamily\small\color{muted}Tanggal produksi} & {\sffamily\small\bfseries 9 Agustus 2026} & {\sffamily\small\color{muted}Tanggal uji} & {\sffamily\small\bfseries 7 September 2026} \\[2pt]
{\sffamily\small\color{muted}Tanggal rilis} & {\sffamily\small\bfseries 20 September 2026} & {\sffamily\small\color{muted}Kedaluwarsa} & {\sffamily\small\bfseries Agustus 2027} \\[2pt]
{\sffamily\small\color{muted}Ukuran batch} & {\sffamily\small\bfseries 100 vial, 5 mL (500 mL)} & {\sffamily\small\color{muted}Penyimpanan} & {\sffamily\small\bfseries 2 sampai 8 $^\circ$C, terlindung cahaya} \\
\end{tabular}

\vspace{3pt}{\sffamily\small\color{muted}Sumber bahan: }{\small Biji gandum, \emph{Triticum aestivum}, tepung tanpa pemutih. Lot sumber GDM-S-07.}\par

\vspace{4pt}
{\footnotesize\tabstart\setlength{\tabcolsep}{3pt}\renewcommand{\arraystretch}{1.08}
\begin{tabular}{@{}rL{3.3cm}L{3.6cm}L{2.9cm}L{2.6cm}L{0.9cm}@{}}
\toprule

\hdrrow\tH{No} & \tH{Parameter} & \tH{Spesifikasi} & \tH{Hasil} & \tH{Metode} & \tH{Status} \\
\midrule

\rowcolor{rule!55}\multicolumn{6}{@{}l}{\sffamily\scriptsize\bfseries\color{tealdark}\MakeUppercase{A. Sifat fisik dan identitas}} \\
1 & Bentuk dan warna & Jernih sampai sedikit opalesen & Jernih sampai sedikit opalesen & Visual & Lulus \\
2 & Volume per vial & 5,0 mL & 5,0 mL & Volumetrik & Lulus \\
3 & Kelarutan & Larut dalam air & Larut sempurna & Visual & Lulus \\
4 & Keasaman (pH) & 6,4 sampai 7,0 & 6,7 & Potensiometri & Lulus \\
\rowcolor{rule!55}\multicolumn{6}{@{}l}{\sffamily\scriptsize\bfseries\color{tealdark}\MakeUppercase{B. Karakterisasi kimia dan biologi}} \\
5 & Nitrogen total & 0,04 sampai 0,09\% & 0,063\% & Kjeldahl & Lulus \\
6 & Protein total & 0,5 sampai 2,5 mg/mL & 1,31 mg/mL & Bradford \cite{ema} & Lulus \\
7 & Profil protein & Pita utama 65 kDa & Sesuai, 65 kDa & SDS-PAGE 12\% & Lulus \\
8 & Kadar Gliadin (Tri a 19) & 50,0 sampai 400,0 ng/mL & 168,0 ng/mL & ELISA gliadin (sandwich) & Lulus \\
9 & Aktivitas alergenik & 80 sampai 125\% IHRP & 93,9\% IHRP & IgE-inhibisi \cite{ema} & Lulus \\
\rowcolor{rule!55}\multicolumn{6}{@{}l}{\sffamily\scriptsize\bfseries\color{tealdark}\MakeUppercase{C. Mutu mikrobiologi dan endotoksin}} \\
10 & Angka lempeng aerob (TAMC) & $\leq 10^2$ CFU/mL & $<10$ CFU/mL & Ph. Eur. 2.6.12 & Lulus \\
11 & Kapang dan khamir (TYMC) & $\leq 10^1$ CFU/mL & $<10$ CFU/mL & Ph. Eur. 2.6.12 & Lulus \\
12 & \emph{Staphylococcus aureus} & Tidak ada per 1 mL & Tidak ada & Ph. Eur. 2.6.13 & Lulus \\
13 & \emph{Pseudomonas aeruginosa} & Tidak ada per 1 mL & Tidak ada & Ph. Eur. 2.6.13 & Lulus \\
14 & Endotoksin & $\leq 20$ EU/mL & 2,4 EU/mL & Ph. Eur. 2.6.14 & Lulus \\

\bottomrule
\end{tabular}}

\vspace{6pt}

\textbf{Kesimpulan.} Batch ALG-GDM-2608-01 memenuhi seluruh spesifikasi dan dirilis sebagai bahan baku patch Alergen.

\vspace{4pt}
\begin{tabular}{@{}p{5.1cm}p{5.1cm}p{5.1cm}@{}}
\rule{0pt}{16pt}\isi[4.5cm] & \isi[4.5cm] & \isi[4.5cm]\\
{\sffamily\scriptsize\color{muted}Analis pengendalian mutu} & {\sffamily\scriptsize\color{muted}Kepala bagian mutu} & {\sffamily\scriptsize\color{muted}Penanggung jawab produksi}\\
{\sffamily\scriptsize\color{muted}Tanggal: 20 September 2026} & {\sffamily\scriptsize\color{muted}Tanggal: 20 September 2026} & {\sffamily\scriptsize\color{muted}Tanggal: 20 September 2026}
\end{tabular}


\clearpage
\subsection{Kacang tanah (\emph{Peanut Extract})}

\coaHeader{Peanut Extract}{Ekstrak alergen bahan baku patch Alergen \ $\cdot$ \ Kacang tanah}

\begin{tabular}{@{}l@{\hspace{8pt}}L{4.3cm}l@{\hspace{8pt}}L{4.3cm}@{}}
{\sffamily\small\color{muted}No. CoA} & {\sffamily\small\bfseries CoA/RM/KCT/08/IX/2026} & {\sffamily\small\color{muted}No. batch} & {\sffamily\small\bfseries ALG-KCT-2608-01} \\[2pt]
{\sffamily\small\color{muted}Tanggal produksi} & {\sffamily\small\bfseries 10 Agustus 2026} & {\sffamily\small\color{muted}Tanggal uji} & {\sffamily\small\bfseries 8 September 2026} \\[2pt]
{\sffamily\small\color{muted}Tanggal rilis} & {\sffamily\small\bfseries 21 September 2026} & {\sffamily\small\color{muted}Kedaluwarsa} & {\sffamily\small\bfseries Agustus 2027} \\[2pt]
{\sffamily\small\color{muted}Ukuran batch} & {\sffamily\small\bfseries 100 vial, 5 mL (500 mL)} & {\sffamily\small\color{muted}Penyimpanan} & {\sffamily\small\bfseries 2 sampai 8 $^\circ$C, terlindung cahaya} \\
\end{tabular}

\vspace{3pt}{\sffamily\small\color{muted}Sumber bahan: }{\small Biji kacang tanah, \emph{Arachis hypogaea}, tanpa kulit ari. Lot sumber KCT-S-08.}\par

\vspace{4pt}
{\footnotesize\tabstart\setlength{\tabcolsep}{3pt}\renewcommand{\arraystretch}{1.08}
\begin{tabular}{@{}rL{3.3cm}L{3.6cm}L{2.9cm}L{2.6cm}L{0.9cm}@{}}
\toprule

\hdrrow\tH{No} & \tH{Parameter} & \tH{Spesifikasi} & \tH{Hasil} & \tH{Metode} & \tH{Status} \\
\midrule

\rowcolor{rule!55}\multicolumn{6}{@{}l}{\sffamily\scriptsize\bfseries\color{tealdark}\MakeUppercase{A. Sifat fisik dan identitas}} \\
1 & Bentuk dan warna & Jernih, kuning pucat & Jernih, kuning pucat & Visual & Lulus \\
2 & Volume per vial & 5,0 mL & 5,0 mL & Volumetrik & Lulus \\
3 & Kelarutan & Larut dalam air & Larut sempurna & Visual & Lulus \\
4 & Keasaman (pH) & 7,0 sampai 7,5 & 7,3 & Potensiometri & Lulus \\
\rowcolor{rule!55}\multicolumn{6}{@{}l}{\sffamily\scriptsize\bfseries\color{tealdark}\MakeUppercase{B. Karakterisasi kimia dan biologi}} \\
5 & Nitrogen total & 0,06 sampai 0,12\% & 0,087\% & Kjeldahl & Lulus \\
6 & Protein total & 0,5 sampai 2,5 mg/mL & 1,88 mg/mL & Bradford \cite{ema} & Lulus \\
7 & Profil protein & Pita utama 17 kDa & Sesuai, 17 kDa & SDS-PAGE 12\% & Lulus \\
8 & Kadar Ara h 2 & 100,0 sampai 800,0 ng/mL & 436,0 ng/mL & ELISA Ara h 2 (sandwich) & Lulus \\
9 & Aktivitas alergenik & 80 sampai 125\% IHRP & 104,5\% IHRP & IgE-inhibisi \cite{ema} & Lulus \\
\rowcolor{rule!55}\multicolumn{6}{@{}l}{\sffamily\scriptsize\bfseries\color{tealdark}\MakeUppercase{C. Mutu mikrobiologi dan endotoksin}} \\
10 & Angka lempeng aerob (TAMC) & $\leq 10^2$ CFU/mL & $<10$ CFU/mL & Ph. Eur. 2.6.12 & Lulus \\
11 & Kapang dan khamir (TYMC) & $\leq 10^1$ CFU/mL & $<10$ CFU/mL & Ph. Eur. 2.6.12 & Lulus \\
12 & \emph{Staphylococcus aureus} & Tidak ada per 1 mL & Tidak ada & Ph. Eur. 2.6.13 & Lulus \\
13 & \emph{Pseudomonas aeruginosa} & Tidak ada per 1 mL & Tidak ada & Ph. Eur. 2.6.13 & Lulus \\
14 & Endotoksin & $\leq 20$ EU/mL & 2,7 EU/mL & Ph. Eur. 2.6.14 & Lulus \\

\bottomrule
\end{tabular}}

\vspace{6pt}

\textbf{Kesimpulan.} Batch ALG-KCT-2608-01 memenuhi seluruh spesifikasi dan dirilis sebagai bahan baku patch Alergen.

\vspace{4pt}
\begin{tabular}{@{}p{5.1cm}p{5.1cm}p{5.1cm}@{}}
\rule{0pt}{16pt}\isi[4.5cm] & \isi[4.5cm] & \isi[4.5cm]\\
{\sffamily\scriptsize\color{muted}Analis pengendalian mutu} & {\sffamily\scriptsize\color{muted}Kepala bagian mutu} & {\sffamily\scriptsize\color{muted}Penanggung jawab produksi}\\
{\sffamily\scriptsize\color{muted}Tanggal: 21 September 2026} & {\sffamily\scriptsize\color{muted}Tanggal: 21 September 2026} & {\sffamily\scriptsize\color{muted}Tanggal: 21 September 2026}
\end{tabular}


\clearpage
\subsection{Kelapa (\emph{Coconut Extract})}

\coaHeader{Coconut Extract}{Ekstrak alergen bahan baku patch Alergen \ $\cdot$ \ Kelapa}

\begin{tabular}{@{}l@{\hspace{8pt}}L{4.3cm}l@{\hspace{8pt}}L{4.3cm}@{}}
{\sffamily\small\color{muted}No. CoA} & {\sffamily\small\bfseries CoA/RM/KLP/09/IX/2026} & {\sffamily\small\color{muted}No. batch} & {\sffamily\small\bfseries ALG-KLP-2608-01} \\[2pt]
{\sffamily\small\color{muted}Tanggal produksi} & {\sffamily\small\bfseries 11 Agustus 2026} & {\sffamily\small\color{muted}Tanggal uji} & {\sffamily\small\bfseries 9 September 2026} \\[2pt]
{\sffamily\small\color{muted}Tanggal rilis} & {\sffamily\small\bfseries 22 September 2026} & {\sffamily\small\color{muted}Kedaluwarsa} & {\sffamily\small\bfseries Agustus 2027} \\[2pt]
{\sffamily\small\color{muted}Ukuran batch} & {\sffamily\small\bfseries 100 vial, 5 mL (500 mL)} & {\sffamily\small\color{muted}Penyimpanan} & {\sffamily\small\bfseries 2 sampai 8 $^\circ$C, terlindung cahaya} \\
\end{tabular}

\vspace{3pt}{\sffamily\small\color{muted}Sumber bahan: }{\small Daging buah kelapa tua, \emph{Cocos nucifera}. Lot sumber KLP-S-09.}\par

\vspace{4pt}
{\footnotesize\tabstart\setlength{\tabcolsep}{3pt}\renewcommand{\arraystretch}{1.08}
\begin{tabular}{@{}rL{3.3cm}L{3.6cm}L{2.9cm}L{2.6cm}L{0.9cm}@{}}
\toprule

\hdrrow\tH{No} & \tH{Parameter} & \tH{Spesifikasi} & \tH{Hasil} & \tH{Metode} & \tH{Status} \\
\midrule

\rowcolor{rule!55}\multicolumn{6}{@{}l}{\sffamily\scriptsize\bfseries\color{tealdark}\MakeUppercase{A. Sifat fisik dan identitas}} \\
1 & Bentuk dan warna & Jernih, tidak berwarna & Jernih, tidak berwarna & Visual & Lulus \\
2 & Volume per vial & 5,0 mL & 5,0 mL & Volumetrik & Lulus \\
3 & Kelarutan & Larut dalam air & Larut sempurna & Visual & Lulus \\
4 & Keasaman (pH) & 6,2 sampai 6,8 & 6,5 & Potensiometri & Lulus \\
\rowcolor{rule!55}\multicolumn{6}{@{}l}{\sffamily\scriptsize\bfseries\color{tealdark}\MakeUppercase{B. Karakterisasi kimia dan biologi}} \\
5 & Nitrogen total & 0,02 sampai 0,05\% & 0,031\% & Kjeldahl & Lulus \\
6 & Protein total & 0,5 sampai 2,5 mg/mL & 0,64 mg/mL & Bradford \cite{ema} & Lulus \\
7 & Profil protein & Pita utama 35 kDa & Sesuai, 35 kDa & SDS-PAGE 12\% & Lulus \\
8 & Kadar Coc n 1 & 40,0 sampai 300,0 ng/mL & 126,0 ng/mL & ELISA Coc n 1 (sandwich) & Lulus \\
9 & Aktivitas alergenik & 80 sampai 125\% IHRP & 96,7\% IHRP & IgE-inhibisi \cite{ema} & Lulus \\
\rowcolor{rule!55}\multicolumn{6}{@{}l}{\sffamily\scriptsize\bfseries\color{tealdark}\MakeUppercase{C. Mutu mikrobiologi dan endotoksin}} \\
10 & Angka lempeng aerob (TAMC) & $\leq 10^2$ CFU/mL & $<10$ CFU/mL & Ph. Eur. 2.6.12 & Lulus \\
11 & Kapang dan khamir (TYMC) & $\leq 10^1$ CFU/mL & $<10$ CFU/mL & Ph. Eur. 2.6.12 & Lulus \\
12 & \emph{Staphylococcus aureus} & Tidak ada per 1 mL & Tidak ada & Ph. Eur. 2.6.13 & Lulus \\
13 & \emph{Pseudomonas aeruginosa} & Tidak ada per 1 mL & Tidak ada & Ph. Eur. 2.6.13 & Lulus \\
14 & Endotoksin & $\leq 20$ EU/mL & 3,0 EU/mL & Ph. Eur. 2.6.14 & Lulus \\

\bottomrule
\end{tabular}}

\vspace{6pt}

\textbf{Kesimpulan.} Batch ALG-KLP-2608-01 memenuhi seluruh spesifikasi dan dirilis sebagai bahan baku patch Alergen.

\vspace{4pt}
\begin{tabular}{@{}p{5.1cm}p{5.1cm}p{5.1cm}@{}}
\rule{0pt}{16pt}\isi[4.5cm] & \isi[4.5cm] & \isi[4.5cm]\\
{\sffamily\scriptsize\color{muted}Analis pengendalian mutu} & {\sffamily\scriptsize\color{muted}Kepala bagian mutu} & {\sffamily\scriptsize\color{muted}Penanggung jawab produksi}\\
{\sffamily\scriptsize\color{muted}Tanggal: 22 September 2026} & {\sffamily\scriptsize\color{muted}Tanggal: 22 September 2026} & {\sffamily\scriptsize\color{muted}Tanggal: 22 September 2026}
\end{tabular}


\clearpage
\subsection{Polen rumput bermuda (\emph{Bermuda Grass Pollen Extract})}

\coaHeader{Bermuda Grass Pollen Extract}{Ekstrak alergen bahan baku patch Alergen \ $\cdot$ \ Polen rumput bermuda}

\begin{tabular}{@{}l@{\hspace{8pt}}L{4.3cm}l@{\hspace{8pt}}L{4.3cm}@{}}
{\sffamily\small\color{muted}No. CoA} & {\sffamily\small\bfseries CoA/RM/PLN/10/IX/2026} & {\sffamily\small\color{muted}No. batch} & {\sffamily\small\bfseries ALG-PLN-2608-01} \\[2pt]
{\sffamily\small\color{muted}Tanggal produksi} & {\sffamily\small\bfseries 12 Agustus 2026} & {\sffamily\small\color{muted}Tanggal uji} & {\sffamily\small\bfseries 10 September 2026} \\[2pt]
{\sffamily\small\color{muted}Tanggal rilis} & {\sffamily\small\bfseries 23 September 2026} & {\sffamily\small\color{muted}Kedaluwarsa} & {\sffamily\small\bfseries Agustus 2027} \\[2pt]
{\sffamily\small\color{muted}Ukuran batch} & {\sffamily\small\bfseries 100 vial, 5 mL (500 mL)} & {\sffamily\small\color{muted}Penyimpanan} & {\sffamily\small\bfseries 2 sampai 8 $^\circ$C, terlindung cahaya} \\
\end{tabular}

\vspace{3pt}{\sffamily\small\color{muted}Sumber bahan: }{\small Polen rumput bermuda, \emph{Cynodon dactylon}, dikumpulkan dan disaring. Lot sumber PLN-S-10.}\par

\vspace{4pt}
{\footnotesize\tabstart\setlength{\tabcolsep}{3pt}\renewcommand{\arraystretch}{1.08}
\begin{tabular}{@{}rL{3.3cm}L{3.6cm}L{2.9cm}L{2.6cm}L{0.9cm}@{}}
\toprule

\hdrrow\tH{No} & \tH{Parameter} & \tH{Spesifikasi} & \tH{Hasil} & \tH{Metode} & \tH{Status} \\
\midrule

\rowcolor{rule!55}\multicolumn{6}{@{}l}{\sffamily\scriptsize\bfseries\color{tealdark}\MakeUppercase{A. Sifat fisik dan identitas}} \\
1 & Bentuk dan warna & Jernih, kuning pucat & Jernih, kuning pucat & Visual & Lulus \\
2 & Volume per vial & 5,0 mL & 5,0 mL & Volumetrik & Lulus \\
3 & Kelarutan & Larut dalam air & Larut sempurna & Visual & Lulus \\
4 & Keasaman (pH) & 7,2 sampai 7,6 & 7,4 & Potensiometri & Lulus \\
\rowcolor{rule!55}\multicolumn{6}{@{}l}{\sffamily\scriptsize\bfseries\color{tealdark}\MakeUppercase{B. Karakterisasi kimia dan biologi}} \\
5 & Nitrogen total & 0,03 sampai 0,08\% & 0,055\% & Kjeldahl & Lulus \\
6 & Protein total & 0,5 sampai 2,5 mg/mL & 1,57 mg/mL & Bradford \cite{ema} & Lulus \\
7 & Profil protein & Pita utama 32 kDa & Sesuai, 32 kDa & SDS-PAGE 12\% & Lulus \\
8 & Kadar Cyn d 1 & 200,0 sampai 1500,0 ng/mL & 812,0 ng/mL & ELISA Cyn d 1 (sandwich) & Lulus \\
9 & Aktivitas alergenik & 80 sampai 125\% IHRP & 100,9\% IHRP & IgE-inhibisi \cite{ema} & Lulus \\
\rowcolor{rule!55}\multicolumn{6}{@{}l}{\sffamily\scriptsize\bfseries\color{tealdark}\MakeUppercase{C. Mutu mikrobiologi dan endotoksin}} \\
10 & Angka lempeng aerob (TAMC) & $\leq 10^2$ CFU/mL & $<10$ CFU/mL & Ph. Eur. 2.6.12 & Lulus \\
11 & Kapang dan khamir (TYMC) & $\leq 10^1$ CFU/mL & $<10$ CFU/mL & Ph. Eur. 2.6.12 & Lulus \\
12 & \emph{Staphylococcus aureus} & Tidak ada per 1 mL & Tidak ada & Ph. Eur. 2.6.13 & Lulus \\
13 & \emph{Pseudomonas aeruginosa} & Tidak ada per 1 mL & Tidak ada & Ph. Eur. 2.6.13 & Lulus \\
14 & Endotoksin & $\leq 20$ EU/mL & 3,3 EU/mL & Ph. Eur. 2.6.14 & Lulus \\

\bottomrule
\end{tabular}}

\vspace{6pt}

\textbf{Kesimpulan.} Batch ALG-PLN-2608-01 memenuhi seluruh spesifikasi dan dirilis sebagai bahan baku patch Alergen.

\vspace{4pt}
\begin{tabular}{@{}p{5.1cm}p{5.1cm}p{5.1cm}@{}}
\rule{0pt}{16pt}\isi[4.5cm] & \isi[4.5cm] & \isi[4.5cm]\\
{\sffamily\scriptsize\color{muted}Analis pengendalian mutu} & {\sffamily\scriptsize\color{muted}Kepala bagian mutu} & {\sffamily\scriptsize\color{muted}Penanggung jawab produksi}\\
{\sffamily\scriptsize\color{muted}Tanggal: 23 September 2026} & {\sffamily\scriptsize\color{muted}Tanggal: 23 September 2026} & {\sffamily\scriptsize\color{muted}Tanggal: 23 September 2026}
\end{tabular}


\clearpage
\section{CoA produk jadi (patch Alergen)}

Hasil per sumur dilampirkan pada Tabel~\ref{tab:sumur}. Stabilitas real-time dipantau terpisah pada 2 sampai 8 $^\circ$C \cite{ich}.

\coaHeader{Patch Alergen, kit diagnostik alergi}{Paten Sederhana IDS000011522 \ $\cdot$ \ 10 ekstrak alergen \ $\cdot$ \ patch 2 $\times$ 5 cm}

\begin{tabular}{@{}l@{\hspace{8pt}}L{4.3cm}l@{\hspace{8pt}}L{4.3cm}@{}}
{\sffamily\small\color{muted}No. CoA} & {\sffamily\small\bfseries CoA/FP/ALG-P/01/X/2026} & {\sffamily\small\color{muted}No. batch} & {\sffamily\small\bfseries ALG-P-2610-01} \\[2pt]
{\sffamily\small\color{muted}Ukuran batch} & {\sffamily\small\bfseries 1.000 patch} & {\sffamily\small\color{muted}Tanggal produksi} & {\sffamily\small\bfseries 6 Oktober 2026} \\[2pt]
{\sffamily\small\color{muted}Tanggal uji} & {\sffamily\small\bfseries 12 Oktober 2026} & {\sffamily\small\color{muted}Tanggal rilis} & {\sffamily\small\bfseries 22 Oktober 2026} \\[2pt]
{\sffamily\small\color{muted}Tempat produksi} & {\sffamily\small\bfseries Fasilitas produksi terkendali, Surabaya} & {\sffamily\small\color{muted}Kedaluwarsa} & {\sffamily\small\bfseries Oktober 2027} \\[2pt]
{\sffamily\small\color{muted}Penyimpanan} & {\sffamily\small\bfseries 2 sampai 8 $^\circ$C, kering, terlindung cahaya} & {\sffamily\small\color{muted}Lot ekstrak} & {\sffamily\small\bfseries 10 batch, Tabel~\ref{tab:ringkasan}} \\
\end{tabular}

\vspace{4pt}
{\scriptsize\tabstart\setlength{\tabcolsep}{3pt}\renewcommand{\arraystretch}{1.08}
\begin{tabular}{@{}rL{3.0cm}L{4.5cm}L{2.9cm}L{2.5cm}L{0.9cm}@{}}
\toprule

\hdrrow\tH{No} & \tH{Parameter} & \tH{Spesifikasi} & \tH{Hasil} & \tH{Metode} & \tH{Status} \\
\midrule

\rowcolor{rule!55}\multicolumn{6}{@{}l}{\sffamily\scriptsize\bfseries\color{tealdark}\MakeUppercase{A. Deskripsi dan dimensi}} \\
1 & Deskripsi & Patch 2 $\times$ 5 cm, 10 sumur, seal utuh & Sesuai & Visual & Lulus \\
2 & Panjang dan lebar & 50 $\times$ 20 mm, toleransi $\pm$ 1 mm & 50,2 $\times$ 20,1 mm & Jangka sorong, n = 10 & Lulus \\
3 & Diameter sumur & 7 mm, toleransi $\pm$ 0,3 mm & 7,0 mm (6,9 sampai 7,1) & Jangka sorong, n = 10 & Lulus \\
4 & Jumlah sumur dan urutan alergen & 10 sumur, urutan sesuai gambar desain & 10 sumur, sesuai & Visual & Lulus \\
\rowcolor{rule!55}\multicolumn{6}{@{}l}{\sffamily\scriptsize\bfseries\color{tealdark}\MakeUppercase{B. Kandungan alergen}} \\
5 & Identitas 10 ekstrak & Profil sesuai IHRP tiap ekstrak \cite{ema} & 10 dari 10 sesuai & SDS-PAGE, ELISA penanda & Lulus \\
6 & Protein per sumur & 30 $\mu$g, batas 24 sampai 36 $\mu$g (80 sampai 120\%) & 28,9 sampai 31,2 $\mu$g & Bradford, n = 10 patch & Lulus \\
7 & Keseragaman antarpatch & RSD $\leq 10\%$ tiap sumur & RSD 3,0 sampai 5,0\% & Dari nomor 6 & Lulus \\
8 & Kadar penanda per sumur & Sesuai rentang tiap ekstrak & 10 dari 10 sesuai & ELISA, n = 10 patch & Lulus \\
9 & Aktivitas alergenik per sumur & 80 sampai 125\% IHRP & 94,6 sampai 103,7\% & IgE-inhibisi \cite{ema} & Lulus \\
10 & Migrasi antarsumur & Tidak terdeteksi (di bawah batas kuantitasi) & Tidak terdeteksi & ELISA silang, n = 3 & Lulus \\

\bottomrule
\end{tabular}}
\vspace{2pt}

\vspace{4pt}
{\scriptsize\tabstart\setlength{\tabcolsep}{3pt}\renewcommand{\arraystretch}{1.08}
\begin{tabular}{@{}rL{3.0cm}L{4.5cm}L{2.9cm}L{2.5cm}L{0.9cm}@{}}
\toprule

\hdrrow\tH{No} & \tH{Parameter} & \tH{Spesifikasi} & \tH{Hasil} & \tH{Metode} & \tH{Status} \\
\midrule

\rowcolor{rule!55}\multicolumn{6}{@{}l}{\sffamily\scriptsize\bfseries\color{tealdark}\MakeUppercase{C. Kemasan dan mutu mikrobiologi}} \\
11 & Integritas seal dan kemasan foil & Tanpa kebocoran atau kerusakan & Sesuai, 0 dari 20 & Uji kebocoran, visual & Lulus \\
12 & Daya rekat & $\geq 1{,}0$ N/cm & 1,6 N/cm & Kupas 180$^\circ$, n = 5 & Lulus \\
13 & Angka lempeng aerob (TAMC) & $\leq 10^2$ CFU per patch \cite{ph514} & $<10$ CFU & Ph. Eur. 2.6.12 & Lulus \\
14 & Kapang dan khamir (TYMC) & $\leq 10^1$ CFU per patch \cite{ph514} & $<10$ CFU & Ph. Eur. 2.6.12 & Lulus \\
15 & \emph{S. aureus} dan \emph{P. aeruginosa} & Tidak ada per patch & Tidak ada & Ph. Eur. 2.6.13 & Lulus \\
16 & Penandaan & Nomor batch, kedaluwarsa, penyimpanan & Sesuai & Visual & Lulus \\

\bottomrule
\end{tabular}}
\vspace{6pt}

\textbf{Kesimpulan.} Batch ALG-P-2610-01 memenuhi seluruh spesifikasi dan dirilis untuk distribusi.

\vspace{4pt}
\begin{tabular}{@{}p{5.1cm}p{5.1cm}p{5.1cm}@{}}
\rule{0pt}{16pt}\isi[4.5cm] & \isi[4.5cm] & \isi[4.5cm]\\
{\sffamily\scriptsize\color{muted}Analis pengendalian mutu} & {\sffamily\scriptsize\color{muted}Kepala bagian mutu} & {\sffamily\scriptsize\color{muted}Penanggung jawab produksi}\\
{\sffamily\scriptsize\color{muted}Tanggal: 22 Oktober 2026} & {\sffamily\scriptsize\color{muted}Tanggal: 22 Oktober 2026} & {\sffamily\scriptsize\color{muted}Tanggal: 22 Oktober 2026}
\end{tabular}


\Needspace{28\baselineskip}
\subsection{Lampiran hasil per sumur}

Tabel~\ref{tab:sumur} merinci protein, penanda, dan aktivitas tiap sumur dari 10 patch yang diuji. Penanda per sumur dihitung dari rasio penanda terhadap protein pada CoA bahan baku, dikalikan protein terukur pada sumur.

\begin{table}[H]
\centering
\caption{Hasil per sumur, batch ALG-P-2610-01 (n = 10 patch).}
\label{tab:sumur}
{\footnotesize\tabstart\setlength{\tabcolsep}{4pt}

\begin{tabular}{@{}rL{2.9cm}L{2.4cm}L{1.1cm}L{2.5cm}L{1.9cm}L{1.1cm}@{}}
\toprule

\hdrrow\tH{Sumur} & \tH{Alergen} & \tH{Protein ($\mu$g)} & \tH{RSD} & \tH{Penanda (ng)} & \tH{Aktivitas} & \tH{Status} \\
\midrule

1 & Tungau debu rumah & 29,4 $\pm$ 1,1 & 3,7\% & 0,49 (Der p 1) & 97,2\% IHRP & Lulus \\

2 & Udang & 30,8 $\pm$ 1,4 & 4,5\% & 9,88 (Tropomiosin) & 102,5\% IHRP & Lulus \\

3 & Susu sapi & 30,1 $\pm$ 0,9 & 3,0\% & 3,89 (Bos d 5) & 99,4\% IHRP & Lulus \\

4 & Daging sapi & 29,7 $\pm$ 1,2 & 4,0\% & 9,96 (Bos d 6) & 96,1\% IHRP & Lulus \\

5 & Daging ayam & 31,2 $\pm$ 1,5 & 4,8\% & 7,14 (Gal d 5) & 101,8\% IHRP & Lulus \\

6 & Putih telur & 28,9 $\pm$ 1,3 & 4,5\% & 2,65 (Ovalbumin) & 94,6\% IHRP & Lulus \\

7 & Gandum & 30,5 $\pm$ 1,0 & 3,3\% & 3,91 (Gliadin) & 100,3\% IHRP & Lulus \\

8 & Kacang tanah & 29,9 $\pm$ 1,1 & 3,7\% & 6,93 (Ara h 2) & 98,9\% IHRP & Lulus \\

9 & Kelapa & 30,3 $\pm$ 1,2 & 4,0\% & 5,97 (Coc n 1) & 103,7\% IHRP & Lulus \\

10 & Polen rumput bermuda & 31,0 $\pm$ 1,4 & 4,5\% & 16,0 (Cyn d 1) & 99,8\% IHRP & Lulus \\

\bottomrule
\end{tabular}}
\end{table}

Batas penerimaan: protein 24 sampai 36 $\mu$g, RSD tidak lebih dari 10\%, aktivitas 80 sampai 125\% IHRP.


\Needspace{30\baselineskip}
\section{Metode uji dan singkatan}

\begin{table}[H]
\centering
\caption{Metode uji yang dipakai pada CoA.}
\label{tab:metode}
{\footnotesize\tabstart\setlength{\tabcolsep}{4pt}
\begin{tabular}{@{}L{3.6cm}L{5.4cm}L{5.6cm}@{}}
\toprule
\hdrrow\tH{Parameter} & \tH{Metode} & \tH{Keterangan} \\
\midrule
Keasaman & Potensiometri dengan pH meter terkalibrasi dua titik & Suhu ruang \\
Nitrogen total & Kjeldahl & Hasil dalam persen bobot per volume \\
Protein total & Bradford, kurva albumin serum sapi & Kit Thermo \#23236 \\
Profil protein & SDS-PAGE 12\%, pewarnaan Coomassie & Dibandingkan dengan IHRP \\
Der p 1 & ELISA & Kit Indoor Biotechnologies \cite{pratama23} \\
Tropomiosin & ELISA & Kit InBio EPC-TPM-1 \cite{marhaeny24} \\
Penanda lain & ELISA sandwich & Kurva standar murni tiap penanda \\
Aktivitas alergenik & Inhibisi pengikatan IgE kompetitif & Serum pool 12 individu atopik, IHRP sebagai acuan \cite{ema} \\
TAMC, TYMC & Ph. Eur. 2.6.12 & Tuang lempeng \\
Mikroorganisme spesifik & Ph. Eur. 2.6.13 & Pengayaan dan isolasi \\
Endotoksin & Ph. Eur. 2.6.14 & Uji LAL kinetik \\
Stabilitas & ICH Q1A(R2), real-time 2 sampai 8 $^\circ$C & Tiap 3 bulan pada tahun pertama \cite{ich} \\
\bottomrule
\end{tabular}}
\end{table}

\paragraph{Singkatan.} CFU, \emph{colony forming unit}; CoA, \emph{certificate of analysis}; ELISA, \emph{enzyme-linked immunosorbent assay}; EU, \emph{endotoxin unit}; IHRP, \emph{in-house reference preparation}; PBS, \emph{phosphate buffered saline}; RSD, simpangan baku relatif; SDS-PAGE, elektroforesis gel poliakrilamida natrium dodesil sulfat; TAMC, \emph{total aerobic microbial count}; TYMC, \emph{total yeasts and moulds count}.

\begin{thebibliography}{99}
\small
\setlength{\itemsep}{2pt}
\bibitem{ema} European Medicines Agency. Guideline on allergen products: production and quality issues. EMEA/CHMP/BWP/304831/2007, 20 November 2008.
\bibitem{ph514} European Pharmacopoeia. Bab 5.1.4, Microbiological quality of non-sterile pharmaceutical preparations and substances for pharmaceutical use; bab 2.6.12, 2.6.13, dan 2.6.14.
\bibitem{ich} International Council for Harmonisation. ICH Q1A(R2): Stability testing of new drug substances and products. 2003.
\bibitem{pratama23} Pratama YA, dkk. The optimisation and standardisation of Indonesian house dust mites allergenic extract as a desensitising agent. Pharmacy Education 2023;23(4):137-139. doi:10.46542/pe.2023.234.137139
\bibitem{marhaeny24} Marhaeny HD, dkk. Shrimp allergen extract immunotherapy induces prolonged immune tolerance in a gastro-food allergy mouse model. PLoS ONE 2024;19(12):e0315312. doi:10.1371/journal.pone.0315312
\end{thebibliography}

\end{document}
